\par\noindent\SourceLeaf{46}%
terminus generalis $r^{\text{tus}}$

\[ = \frac{1}{2^{2m-1}}{}^{2m}\!\mathfrak{B}^{(m+r)}\! . \cos 2r\varphi. \]

$\beta$. Potestates impares eidem legi subiectae inveniuntur, nisi quod
\noindent terminus initialis in numeratore, coefficientium binomialium, ad
potestatem datam pertinentium medium posteriorem (quippe bini in exponente
impari occurrunt et medii et aequales) totum contineat. Lex seriei
generalis igitur haec erit: esse $\cos\varphi^{2m-1}$ seriem, cuius
\noindent terminus initialis

\[ = \frac{1}{2^{2m-2}}{}^{2m-1}\!\mathfrak{B}^{(m)}\cos\varphi, \]

\noindent terminus generalis $r^{\text{tus}}$

\[ = \frac{1}{2^{2m-2}}{}^{2m-1}\!\mathfrak{B}^{(m+r)}\cos(2r+1)\varphi. \]

3. Sumamus igitur primo loco legem propositam in indefinita aliqua
potestate impari, v.\,g.\ $(2m-1)^{\text{ta}}$ valere, adeoque esse

\[
\begin{aligned}
\cos\varphi^{2m-1} & = \frac{1}{2^{2m-2}}{}^{2m-1}\!\mathfrak{B}^{(m)}\cos\varphi +\\
&{} \frac{1}{2^{2m-2}}{}^{2m-1}\!\mathfrak{B}^{(m+1)}\cos 3\varphi +\\
&{} \frac{1}{2^{2m-2}}{}^{2m-1}\!\mathfrak{B}^{(m+2)}\cos 5\varphi + \text{etc.},
\end{aligned}
\]

\noindent obtinebitur potestas par unitate maior, utramque aequationis partem per
$\cos\varphi$ multiplicando. Quodlibet serie membrum hoc pacto productum
binorum cosinuum continet, nec difficulter ex theoremate noto in summam
binorum cosinuum solvetur. Sic fiet

\[
\begin{aligned}
\cos\varphi^{2m} & = \frac{1}{2^{2m-2}}{}^{2m-1}\!\mathfrak{B}^{(m)}\cos\varphi . \cos\varphi +\\
&{} \frac{1}{2^{2m-2}}{}^{2m-1}\!\mathfrak{B}^{(m+1)}\cos 3\varphi . \cos\varphi +\\
&{} \frac{1}{2^{2m-2}}{}^{2m-1}\!\mathfrak{B}^{(m+2)}\cos 5\varphi . \cos\varphi + .\,.\,.
\end{aligned}
\]

\noindent Evadet igitur, quia

\[ \cos\varphi\cos\varphi = \tfrac12+\tfrac12\cos 2\varphi, \]

\noindent terminus initialis

\[ = \frac{1}{2^{2m-2}}{}^{2m-1}\!\mathfrak{B}^{(m)} . \tfrac12 +
\frac{1}{2^{2m-2}}{}^{2m-1}\!\mathfrak{B}^{(m)}\frac{\cos 2\varphi}{2}, \]

\noindent eademque lege terminus primus

\[ = \frac{1}{2^{2m-2}}{}^{2m-1}\!\mathfrak{B}^{(m+1)} . \frac{\cos 2\varphi}{2} +
\frac{1}{2^{2m-2}}{}^{2m-1}\!\mathfrak{B}^{(m+1)}\frac{\cos 4\varphi}{2}. \]
